\section{Discusión}\label{sec:discussion}
A3 demuestra reducción de parámetros y una ventaja observada de memoria frente a E5 sobre los servicios nuevos fijos. La evidencia completa no respalda mejora uniforme. CPU A3 empeora frente a ambas referencias principales, A0 tiene menor memoria MASE que A3 y PatchTST pequeño obtiene la menor media principal. La ventaja contra E5 se concentra en una minoría de servicios, aunque la dirección por semilla es consistente bajo la comparación preespecificada. Las exclusiones posteriores informadas por resultados explican por qué esa consistencia no equivale a portabilidad espacial amplia.

El resultado algebraico delimita la aportación: dos cabezas afines tras un filtro móvil fijo siguen siendo una función afín. Reducir parámetros puede bajar el coste de ejecución medido, pero cambia la trayectoria de Adam y la interacción con regularización. Los coeficientes de ausencia pueden mejorar o empeorar. EMA dependiente de disponibilidad en E5 no queda incluida en ese colapso. Las medias A1/A2/A3 seleccionadas independientemente no prueban un beneficio aditivo general. Un estudio mecanístico posterior podría fijar hiperparámetros, igualar regularización y evaluar interacciones sobre otro período independiente; son propuestas, no experimentos ejecutados.

La transición histórica-ampliada modifica más que las semillas. Los servicios destino tienen distinta calibración, otro papel en el ajuste y una receta fuente con diez condiciones. El cambio de rango de E5 no puede atribuirse únicamente a aumentar servicios o inicializaciones. La mejora histórica de memoria, su conservación incompleta por C/V2 y el contraste A3 posterior responden conjuntamente a la parsimonia: reducir coeficientes puede mantener un error observado aceptable y perder una ventaja previa específica. Una comparación favorable posterior tampoco restablece superioridad universal.

NLinear y PatchTST sugieren competitividad de un nivel reciente explícito y una representación por parches bajo transferencia fuente. No se aíslan esos mecanismos: cambian configuraciones, parametrización y, en PatchTST, atención y normalización. NLinear tiene MAE CPU mostrado ligeramente menor que PatchTST pese al orden contrario de MASE. MASE pondera respecto de variación histórica por servicio; el error absoluto conserva la unidad física. La clasificación depende de la función de pérdida y la cohorte, sin que una métrica invalide automáticamente a la otra.

Los horizontes directos mayores aumentan errores mostrados de los modelos seleccionados, pero utilizan pesos propios y objetivos comunes; no prueban acumulación recursiva. Las historias largas producen mejora CPU y empeoramiento memoria en A3. Contener más valores aumenta coeficientes, procesamiento y exposición potencial a regímenes antiguos; no garantiza generalización. El diseño estrella describe elecciones concretas, sin estimar todas las interacciones ventana-horizonte.

El caso externo cambia dominio y definición de medición. La ventaja de persistencia es compatible con una traza corta de utilización normalizada relativamente estable, pero no establece su causa. No hubo grafo externo y seis horas pueden producir denominadores pequeños. Se muestran MASE y errores absolutos normalizados juntos, sin combinar magnitudes con errores físicos ChronoGraph. Más días, servicios elegidos independientemente y datasets con preparación documentada permitirían estudiar estabilidad de rangos.

La eficiencia ofrece una perspectiva acotada. A3 es menor que E5 y tiene menor p 50 con tensores residentes en ambos dispositivos. PatchTST presenta menor MASE principal con mayor coste medido; media móvil resulta una referencia CPU fuerte sin pesos entrenables. En lotes pequeños, lanzamiento y sincronización pueden favorecer CPU frente a CUDA. No se demuestra menor energía, recolección más rápida o mejor autoescalado. Evaluarlo exigiría controlador, reproducción de cargas, tiempos completos, métricas SLA e instrumentación energética atribuible al proyecto.

La corrección histórica de unidades de intervalos es sustantiva: un intervalo absoluto MASE no puede leerse como porcentaje porque aparezca junto a una reducción porcentual. En la ampliación, porcentajes por celda y cociente de medias se mantienen separados. La prueba de signos mide consistencia de dirección entre ajustes y el intervalo por bloques aproxima una media remuestreada. Su desacuerdo es posible sin error aritmético y no permite sustituir el criterio principal por el más favorable. Así se evita que un rango o una incertidumbre condicional estrecha se convierta en superioridad poblacional no demostrada.

El histórico conserva alcance distinto y genera hipótesis. La transferencia espacial cambia el orden y Alibaba lo cambia otra vez. El extremo GinAR se retiene como resultado de implementación/protocolo, sin refutar el método publicado. La contribución reproducible es documentar cuándo capacidad, procedencia de observaciones, error, coste y dominio dejan de coincidir en una misma recomendación.

La dinámica calculada solo en entrenamiento aporta contexto para interpretar los dos recursos. La autocorrelación mediana de primer retardo entre las 28 series es 0.447 en CPU y 0.867 en memoria. El cambio absoluto relativo a la desviación de cada serie tiene medianas 0.413 y 0.275. Son valores compatibles con mayor persistencia de nivel en memoria. No identifican la causa del cambio de error con EMA: el residual conserva las desviaciones del nivel y ambas cabezas se ajustan juntas. Motivan un análisis controlado alrededor de cambios abruptos, sin afirmar que suavizar retrasa mecánicamente el pronóstico final.
